% GENERATED from QA2-the-skill-set.md by md2tex.py. Do not hand-edit: sync is
% one-way and the next run overwrites this file.

\section{📦 The roster: 23 skills and 8 packs}
\textbf{Counts by layer}: the family tree as it sits on disk today.

\begin{verbatim}
  door         1 skill    the router and console front door
 0 0-enter      2 skills   enter . round
 1 1-lifecycle  11 skills  1 orchestrator + 10 per-stage
 2 2-phase      4 skills   draft . probe . revise . check
 3 3-deliver    4 skills   artifact . review . claim-audit . deploy
 4 4-iterate    1 skill    post-deploy A/B refinement
  venue/       8 packs    knowledge, not stages
\end{verbatim}
📌 This part is the count itself: every callable skill with its version, plus the venue packs beside them, read from disk on 260802. The layer grouping follows the Skill-tree layout in the family README, which is the canonical structure file.

\textbf{The roster, grouped by layer}: ver is the skill's own \texttt{metadata.version}, updated is its \texttt{last\_updated}.

\begin{verbatim}
 layer         skill                              ver      updated
 -----------   --------------------------------   ------   ----------
  door       haipipe-application                0.6.10   2026-07-19
 0-enter       haipipe-application-enter          0.2.3    2026-07-19
               haipipe-application-round          0.1.1    2026-07-06
 1-lifecycle   haipipe-application-lifecycle      0.4.4    2026-07-19
    stage    haipipe-application-seed           0.6.1    2026-07-19
    stage    haipipe-application-descriptions   0.2.6    2026-07-19
    stage    haipipe-application-themes         0.2.7    2026-07-19
    stage    haipipe-application-claims         0.7.6    2026-07-19
    stage    haipipe-application-advice         0.1.9    2026-07-19
    stage    haipipe-application-venue          0.3.4    2026-07-19
    stage    haipipe-application-pitch          0.5.4    2026-07-19
    stage    haipipe-application-narrative      0.5.4    2026-07-19
    stage    haipipe-application-display        0.4.8    2026-07-19
    stage    haipipe-application-section-edit   0.5.4    2026-07-19
 2-phase       haipipe-application-draft          0.1.5    2026-07-19
               haipipe-application-probe          0.3.2    2026-07-19
               haipipe-application-revise         0.1.1    2026-07-19
               haipipe-application-check          0.4.2    2026-07-17
 3-deliver     haipipe-application-artifact       0.3.1    2026-07-17
               haipipe-application-review         0.1.1    2026-07-06
               haipipe-application-claim-audit    0.1.2    2026-07-17
               haipipe-application-deploy         0.1.1    2026-07-06
 4-iterate     haipipe-application-iterate        0.2.0    2026-07-17
 -----------   --------------------------------   ------   ----------
               23 skills .  ×10 = round-3 collapse candidate
\end{verbatim}
The ten 🔻 rows are the per-stage skills of \texttt{1-lifecycle/}: the nine numbered stage folders plus \texttt{haipipe-application-venue}, which pins the modality between the ladder and pitch. The orchestrator \texttt{haipipe-application-lifecycle} is the eleventh skill in that layer and is not a candidate; it is the seat the engine proposal builds beside.

Under \texttt{venue/} sit eight packs: sms, email, push, reminder, checklist, dashboard, report, and ui-card. Each is a README plus a style-profile, with a shared \texttt{\_SCHEMA.md} beside them, and the family README calls them knowledge, not stages. They carry no version and no trigger, so they ship with the family without being part of the 23.

Nothing has reached 1.0: the highest version anywhere is \texttt{haipipe-application-claims} at 0.7.6, and seven skills are still 0.1.x. Sixteen of the 23 last moved on 2026-07-19 and nothing has moved since, so this count is stable rather than mid-flight. The naming is uniform: all 23 carry the \texttt{haipipe-application-} prefix, which QA2@paper could not say of its family, where four skills dropped the prefix. The ladder's wear is uneven: \texttt{haipipe-application-claims} sits at 0.7.6 while \texttt{haipipe-application-advice} sits at 0.1.9, because the 1d rung was renamed from \texttt{1d-principles} to \texttt{1d-advice} only on 2026-07-09 and its skill is weeks old.


\section{🔻 The round-3 collapse}
\textbf{Before and after the collapse}: the shapes option A trades between.

\begin{verbatim}
  today      1-lifecycle/ = 1 orchestrator + 10 per-stage SKILL.md
 [x] option A   1-lifecycle/ = orchestrator + 1 engine + stage.md ×10
  precedent  paper: haipipe-paper-stage + 8 stage contracts (QA2@paper)
\end{verbatim}
📌 This part names what round 3 would move, the precedent it copies, and the files a yes must land in.

The ten candidates are all of \texttt{1-lifecycle/} except the orchestrator: seed, descriptions, themes, claims, advice, venue, pitch, narrative, display, and section-edit. Each would stop shipping a SKILL.md and become one \texttt{stages/<dir>/stage.md} data file under a new \texttt{haipipe-application-stage} engine. The door, the 0-enter pair, the four phase workers, the four deliver skills, iterate, and the lifecycle orchestrator stay callable, so the family lands at 14 skills. The venue packs are untouched: they are knowledge the aligned stages consult, whichever form the stages take.

The paper family took this shape in its 260719 refactor: per-stage skills became stage contracts under one runner. QA2@paper's roster shows the result, \texttt{haipipe-paper-stage} at 0.6.0, and its Files row carries \texttt{stages/index.yml}, the eight contracts, and \texttt{CONTRACT.md}. The application README declares the two families structural twins, same spine, same phases, same probe door, so a shape proven on one side is the default proposal for the other.

QA2@paper states the closing rule this page inherits: a ruling that reaches [x] has to land somewhere concrete or it has not graduated. For option A the landing list is: a new \texttt{1-lifecycle/haipipe-application-stage/SKILL.md}, ten \texttt{stages/<dir>/stage.md} files, the ten retired skill folders removed from the tree, the README's Skill-tree layout rewritten, and the door's router table repointed. A ruling recorded only as architecture prose on this board would be exactly the drift that rule exists to stop.

JL selected a Page-first option that the original A/B/C row did not contain. The canonical files are \texttt{haipipe-application/SKILL.md} plus \texttt{page-types/haipipe-page-for-\{brief,intervention,artifact\}/SKILL.md}. The door folds legacy Seed/Venue/Pitch into Brief, Narrative/Display into Intervention, and Section Edit/Draft into Artifact. Descriptions/Themes/Claims/Advice do not become Application stages: missing knowledge routes to \texttt{/haipipe-task insight}.

